\chapter{Introduccion y objetivos del proyecto}\label{defobjetivos}

\section{Contexto y motivacion}

El consumidor espa\~nol afronta la compra cotidiana en un contexto de fuerte variabilidad de
precios entre cadenas, ausencia de una fuente unificada y alto coste temporal para comparar
ofertas manualmente. Aunque el ahorro potencial entre supermercados es significativo, en la
practica ese ahorro se pierde cuando no se considera el coste de desplazamiento ni el tiempo de
ruta.

BargAIn nace para resolver ese problema como una plataforma de compra inteligente que integra
precio, distancia y tiempo en una unica decision. El sistema combina una app movil para
consumidores, un portal web para PYMEs, y un backend modular con capacidades geoespaciales,
OCR y asistencia conversacional.

\section{Problema abordado}

Las soluciones actuales de comparacion de supermercados suelen responder a la pregunta
``que tienda tiene el precio mas bajo'', pero no a la pregunta completa ``que compra me conviene
realmente segun mi ubicacion y mis preferencias''.

El problema de BargAIn se formula como una optimizacion multicriterio en la que intervienen:

\begin{itemize}
	\item coste total de la cesta,
	\item distancia adicional de desplazamiento,
	\item tiempo total estimado de compra y ruta.
\end{itemize}

\section{Objetivo principal}

Desarrollar una aplicacion movil y web que optimice la cesta de la compra del usuario,
calculando la combinacion de supermercados y comercios locales que ofrece la mejor relacion
precio--distancia--tiempo segun preferencias configurables.

\section{Objetivos especificos}

\begin{enumerate}
	\item Implementar una API REST con Django y DRF para usuarios, productos, tiendas,
		precios, listas, optimizacion, OCR, asistente y portal business.
	\item Integrar una capa geoespacial con PostgreSQL + PostGIS para consultas de proximidad
		y preparacion de rutas.
	\item Desarrollar un optimizador multicriterio con OR-Tools y una funcion de scoring
		configurable por el usuario.
	\item Incorporar OCR con Google Cloud Vision para digitalizar tickets y listas,
		con validacion y correccion manual en frontend.
	\item Integrar un asistente LLM (Gemini via backend proxy) con guardrails de dominio
		y contexto de compra.
	\item Construir un portal para PYMEs con gestion de precios y promociones.
	\item Asegurar calidad con tests unitarios, de integracion y E2E, ademas de UAT
		manual para flujos nativos.
	\item Preparar despliegue reproducible en staging con Render, Docker Compose y CI.
\end{enumerate}

\section{Alcance del proyecto}

El alcance del TFG cubre el ciclo completo de ingenieria del software:

\begin{itemize}
	\item analisis y requisitos,
	\item diseno arquitectonico y de datos,
	\item implementacion backend y frontend,
	\item pruebas funcionales y no funcionales,
	\item despliegue en entorno staging,
	\item documentacion tecnica y memoria final.
\end{itemize}

Quedan fuera del alcance de esta version: despliegue productivo global,
acuerdos formales con todas las cadenas para APIs oficiales, y la version final
de Live Activities nativas en iOS (ADR-010).

\section{Tecnologias clave}

El sistema se apoya en un stack full-stack moderno y coherente:

\begin{itemize}
	\item \textbf{Backend:} Django 5, DRF, Celery, Redis.
	\item \textbf{Datos:} PostgreSQL 16 + PostGIS 3.4.
	\item \textbf{Frontend:} React Native con Expo + companion web en React.
	\item \textbf{IA:} Google Cloud Vision API (OCR) y Gemini via backend proxy.
	\item \textbf{Optimizacion:} OR-Tools + calculo de distancias (ORS + fallback haversine).
	\item \textbf{Ops:} Docker, GitHub Actions, Render, Sentry.
\end{itemize}

\section{Estructura de la memoria}

La memoria se organiza para reflejar el recorrido completo del proyecto: contexto y objetivos,
estado del arte, comparativa, planificacion, requisitos, diseno e implementacion,
manual de usuario, pruebas y resultados, conclusiones y bibliografia.
